\SourceLeaf{231}%
Ich habe hier bloss die Grundzüge des Beweises angeben wollen, ohne alle Feile oder Politur, die ich ihm zu geben jetzt keine Zeit habe. Es steht Dir frei, es Deinem Sohne mitzutheilen: jedenfalls bitte ich Dich, ihn herzlich von mir zu grüssen und ihm meine besondere Hochachtung zu versichern; fordere ihn aber doch zugleich auf, sich mit der Aufgabe zu beschäftigen:

»Den Kubikinhalt des Tetraeders (von vier Ebenen begrenzten Raumes) zu bestimmen«.

Da der Flächeninhalt eines Dreiecks sich so einfach angeben lässt: so hätte man erwarten sollen, dass es auch für diesen Kubikinhalt einen eben so einfachen Ausdruck geben werde: aber diese Erwartung wird, wie es scheint, getäuscht.

Um die Geometrie vom Anfange an ordentlich zu behandeln, ist es unerlässlich, die Möglichkeit eines Planums zu beweisen; die gewöhnliche Definition enthält zu viel, und implicirt eigentlich subreptive schon ein Theorem. Man muss sich wundern, dass alle Schriftsteller von Euklid bis auf die neuesten Zeiten so nachlässig dabei zu Werk gegangen sind: allein diese Schwierigkeit ist von durchaus verschiedener [Natur] mit der Schwierigkeit zwischen $\Sigma$ und $S$ zu entscheiden, und jene ist nicht gar schwer zu heben. Wahrscheinlich finde ich mich auch schon durch Dein Buch hierüber befriedigt.

Gerade in der Unmöglichkeit, zwischen $\Sigma$ und $S$ a priori zu entscheiden, liegt der klarste Beweis, dass Kant Unrecht hatte zu behaupten, der Raum sei nur \textls{Form} unserer Anschauung. Einen andern ebenso starken Grund habe ich in einem kleinen Aufsatze angedeutet, der in den Göttingischen Gelehrten Anzeigen 1831 steht Stück 64, pag.~625. Vielleicht wird es Dich nicht gereuen, wenn Du Dich bemühest Dir diesen Band der G.G.A. zu verschaffen (was jeder Buchhändler in Wien oder Ofen leicht bewirken kann), da darin unter andern auch die Quintessenz meiner Ansicht von den imaginären Grössen auf ein Paar Seiten dargelegt ist. \mbox{.\ .\ .\ .\ .\ .\ .}

